\documentclass[11pt,a4paper]{article}
\usepackage[a4paper,margin=1in]{geometry}
\usepackage[utf8]{inputenc}
\usepackage[T1]{fontenc}
\usepackage{microtype}
\usepackage{booktabs}
\usepackage{xcolor}
\usepackage{listings}
\usepackage{enumitem}
\usepackage{hyperref}
\usepackage{tikz}
\hypersetup{colorlinks=false, hidelinks}

\lstset{
  basicstyle=\ttfamily\footnotesize,
  breaklines=true,
  columns=fullflexible,
  keepspaces=true,
  frame=none,
  xleftmargin=1em,
}

\title{The Shakedown\\
\large A Companion to the Entmoot Architecture Paper}
\author{}
\date{April 2026}

\begin{document}
\maketitle

\begin{abstract}
\noindent
Between 17 and 27 April 2026 we shipped a dense sequence of
patch releases across two layered systems: the Pilot fork moved from
\texttt{v1.7.2-jf.1} to \texttt{v1.9.0-jf.15.21}, and Entmoot moved
from \texttt{v1.0.0} to \texttt{v1.5.10}. Every patch was
driven by a symptom observed on a three-machine live demo across
three continents: a laptop on a UAE carrier-grade NAT, a Raspberry Pi
in Italy on Fastweb residential, and a Hetzner VPS in Germany. The
CHANGELOGs document what changed, but their
language assumes a reader who already speaks networking. This
document is the narrative companion: it explains what the words
mean, what actually went wrong on the wire, and why each fix is
the canonical one. It is meant for anyone who wants to understand
the project without first reading RFCs.
\end{abstract}

\tableofcontents
\vspace{1em}

\section{What this document is}
\label{sec:intro}

The Entmoot research paper (\texttt{paper/main.tex}) describes the
protocol as it is designed. This companion describes what went
wrong in the first field deployment and how each problem was
resolved. If the research paper is the blueprint, this is the
engineer's diary kept while wiring the building.

Two things distinguish this document from the CHANGELOGs:

\begin{itemize}[leftmargin=1.2em]
    \item It assumes you have never heard of NAT traversal,
    gossip protocols, or Merkle trees. Section~\ref{sec:mental}
    sets up the mental model; Section~\ref{sec:glossary} defines
    every technical term used later.
    \item It explains each fix in four parts: what the user saw,
    what was really happening, what we changed, and why that
    change is the right shape. Each piece of jargon is introduced
    in context the first time it appears.
\end{itemize}

If you only read three sections, read Section~\ref{sec:mental} for
the mental model, Section~\ref{sec:glossary} for vocabulary, and
Section~\ref{sec:bigpicture} for the shape of what went wrong.

\section{The world before the shakedown}
\label{sec:mental}

\subsection{Three machines, three layers}

The deployment under test has three participants:

\begin{description}[leftmargin=1.2em]
    \item[Laptop] A MacBook in the United Arab Emirates, behind a
    mobile carrier's network. The carrier's equipment rewrites
    every outbound packet so the laptop has no meaningful
    ``address'' on the public internet.
    \item[Phobos] A Raspberry Pi sitting on a residential Italian
    internet connection. Also behind a home router that rewrites
    addresses.
    \item[VPS] A small Linux server rented from Hetzner in
    Germany. It has a real public IP address that the internet
    can reach directly.
\end{description}

Each machine runs two programs that work together:

\begin{description}[leftmargin=1.2em]
    \item[\texttt{pilot-daemon}] The lower layer. Its job is to
    open an encrypted pipe to any other Pilot node, identified by
    a stable numeric address, regardless of which network the
    other node sits on. Think of it as the post office that
    figures out the physical route from sender to recipient.
    \item[\texttt{entmootd join}] The upper layer. Its job is to
    participate in a group: receive messages that anyone else in
    the group publishes, and push its own. Think of it as the
    language the letters are written in, on top of the postal
    service underneath.
\end{description}

The line between the two matters. Many of the fixes in this
document look like problems in the upper layer but were really
problems in the lower layer percolating upward. See
Section~\ref{sec:bigpicture}.

\subsection{Five metaphors that are load-bearing}

These stand in for full technical definitions throughout the
document. The precise definitions are in Section~\ref{sec:glossary}.

\begin{description}[leftmargin=1.2em]
    \item[NAT is a hotel switchboard.] When you dial out from
    your hotel room, the switchboard remembers the outbound line
    and routes any reply back to your room. If the switchboard
    has no prior record of you calling the outside party, an
    incoming call cannot reach you. Home routers, mobile
    carriers, and corporate firewalls all act as switchboards.
    NAT traversal is the art of getting two switchboards to each
    open a line at the same moment so voices can flow.
    \item[Gossip is rumour spreading.] When a node learns
    something new, it whispers the news to a small random subset
    of its neighbours. Each of them whispers onward until
    everyone knows. The mathematics works out: with enough
    whispering, convergence is fast and robust even if some
    neighbours are offline.
    \item[A Merkle root is an inventory total.] Two warehouses
    can confirm their stock is identical by comparing a single
    number (the total), rather than listing every box. If the
    totals match, the inventories are identical; if they differ,
    the warehouses drill into their records to find the gap.
    Merkle trees make this check cryptographically strong.
    \item[A handshake is a mutual nameplate exchange.] Neither
    side speaks in earnest until both have shown signed
    identification and exchanged a shared secret. Afterwards
    their conversation is encrypted with keys only they can
    derive.
    \item[A relay is a forwarding post office.] When two hosts
    cannot reach each other directly, a mutual acquaintance
    passes letters between them. Slower than direct delivery,
    but the only option when both sides are behind
    switchboards that will not cooperate.
\end{description}

\subsection{The four verbs you will see}

\begin{description}[leftmargin=1.2em]
    \item[Join.] A machine loads an invite (a signed document
    that names the group and the current roster) and starts
    participating. Joining involves contacting at least one
    existing member, downloading missing history, and announcing
    itself.
    \item[Publish.] Author a message, sign it, store it locally,
    hand it to the gossip layer for dissemination.
    \item[Tail.] Subscribe to live messages in a group.
    \item[Query.] Read the local store directly for
    already-received messages, with filters for author, topic,
    and time.
\end{description}

\section{Glossary}
\label{sec:glossary}

Every term bolded on first use. References point at the
authoritative source for each concept.

\subsection*{Cryptography and identity}

\begin{description}[leftmargin=1em,style=nextline]
    \item[\textbf{Ed25519 signature.}]
    A fast, deterministic digital signature based on elliptic
    curves. Used to prove a message was written by a specific
    key. (RFC 8032.)
    \item[\textbf{X25519 key exchange.}]
    A procedure that lets two parties derive the same secret
    without anyone listening being able to learn it. (RFC 7748.)
    \item[\textbf{AES-256-GCM.}]
    A symmetric encryption scheme that also detects tampering,
    used to encrypt messages once X25519 has produced a shared
    key. (NIST SP 800-38D.)
    \item[\textbf{Signed append-only log.}]
    A list of entries that can only grow. Each entry is signed,
    so nothing can be silently edited after the fact. Entmoot's
    group roster is one of these. (Certificate Transparency,
    RFC 9162, uses the same primitive.)
    \item[\textbf{Handshake / trust pair.}]
    A one-time mutual authorisation step where two Pilot daemons
    each approve the other before traffic flows. Without this
    pair the daemon will reject the other's packets at the front
    door.
    \item[\textbf{Rekey.}]
    Replacing session keys periodically, so a leaked key only
    compromises a short window of past traffic. (NIST SP 800-57.)
\end{description}

\subsection*{Gossip and dissemination}

\begin{description}[leftmargin=1em,style=nextline]
    \item[\textbf{Fanout.}]
    How many peers a publisher pushes a new message to. Each of
    those peers may fanout again, so a small fanout per hop
    still reaches everyone quickly.
    \item[\textbf{Plumtree (eager + lazy push).}]
    A gossip design that builds a spanning tree for fast
    delivery of full message bodies (eager push) and falls back
    to broadcasting only message IDs on non-tree links (lazy
    push), to repair gaps cheaply. (Leit\~ao et al., SRDS 2007.)
    \item[\textbf{IHAVE / GRAFT / PRUNE.}]
    The three control messages that make Plumtree's tree
    self-healing. IHAVE says ``I have message X but am not
    sending it to you''; GRAFT says ``I did not receive X
    within the timeout, please send it and promote me back onto
    the eager branch''; PRUNE says ``I already had X via
    another path, stop sending duplicates, demote me off the
    eager branch''.
    \item[\textbf{First-seen forwarding.}]
    The rule that every node forwards a message exactly the
    first time it sees it, regardless of which path it arrived
    on. Duplicates are dropped. (Demers et al., PODC 1987.)
    \item[\textbf{Anti-entropy reconciliation.}]
    A background synchronisation round between two peers that
    compares inventory (Merkle roots), finds the gap, and
    exchanges the missing messages. Runs when a peer reconnects
    after a disconnection, and periodically otherwise.
    Scuttlebutt, Riak, and Cassandra all use variants.
    \item[\textbf{Retry budget.}]
    The maximum number of times a transient failure will be
    retried before giving up. Bounded so that one dead peer
    cannot cause infinite work. (Google SRE Book, ``Handling
    Overload''.)
    \item[\textbf{Decorrelated-jitter backoff.}]
    A retry-delay formula that picks the next wait randomly
    between a base and three times the previous wait. Unlike
    fixed exponential doubling, this prevents many nodes from
    retrying in lockstep after a shared hiccup. (Marc Brooker,
    AWS 2015.)
    \item[\textbf{Dial-backoff cache.}]
    A per-peer memory of recent failed connection attempts.
    Further attempts to the same address short-circuit
    immediately until the window reopens. (go-libp2p
    \texttt{swarm.DialBackoff}.)
    \item[\textbf{MQTT topics.}]
    Slash-separated names for subjects (\texttt{chat/eng/alerts}).
    Subscribers can use \texttt{+} to match any one segment and
    \texttt{\#} to match the tail. (MQTT v5 OASIS spec.)
\end{description}

\subsection*{Addresses, NAT, and relays}

\begin{description}[leftmargin=1em,style=nextline]
    \item[\textbf{NAT (Network Address Translation).}]
    A router function that maps many internal addresses to one
    public one. (RFC 3022.)
    \item[\textbf{RFC 1918 address.}]
    An internal-only IPv4 range that is never routable on the
    public internet: \texttt{10.0.0.0/8}, \texttt{172.16.0.0/12},
    \texttt{192.168.0.0/16}. If you see one in a network trace
    it is a sign that the real address has been obscured.
    \item[\textbf{Carrier-grade NAT (CGN / CGNAT).}]
    A second layer of NAT applied by the ISP, sharing a small
    pool of public addresses across many subscribers. Common in
    mobile and recent residential deployments. (RFC 6598.)
    \item[\textbf{STUN.}]
    A short protocol a host can run against a well-known server
    to learn what public address its NAT currently presents to
    the outside world. (RFC 8489.)
    \item[\textbf{NAT hole-punching.}]
    Two hosts behind NATs send packets toward each other at the
    same moment so that both NATs open a mapping on each side,
    letting subsequent traffic flow through. (RFC 5128, \S3.3.)
    \item[\textbf{Symmetric NAT.}]
    A NAT that chooses a new external port per destination,
    which breaks naive hole-punching because the port you
    learned via STUN is not the port another peer will dial.
    (RFC 3489, \S5.)
    \item[\textbf{Relay fallback.}]
    When two peers cannot reach each other directly, a third
    host reachable to both forwards packets between them.
    (TURN, RFC 8656.)
    \item[\textbf{Registry / rendezvous server.}]
    The well-known server that every Pilot daemon contacts on
    startup to announce its current address and look up peers.
    Pilot uses \texttt{34.71.57.205}.
    \item[\textbf{Beacon.}]
    A companion to the registry that handles NAT-punch
    coordination and relays traffic between peers that need one.
    \item[\textbf{Keepalive.}]
    A small periodic packet sent on an otherwise-idle tunnel so
    the NAT does not forget the mapping. (RFC 1122, \S4.2.3.6.)
    \item[\textbf{Spurious RST.}]
    A reset packet that tears down a TCP-like connection without
    a genuine protocol reason, usually caused by a middlebox or
    stale NAT state. (RFC 5961.)
\end{description}

\subsection*{Pilot-specific terms}

\begin{description}[leftmargin=1em,style=nextline]
    \item[\textbf{Node ID.}]
    A 32-bit integer Pilot assigns at registration, bound to the
    daemon's Ed25519 key. The group uses these as stable
    addresses.
    \item[\textbf{PILA.}]
    Pilot's authenticated handshake and key-exchange frame. When
    a peer cannot decrypt a packet, sending an unsolicited PILA
    asks the other side to re-handshake.
    \item[\textbf{Tunnel.}]
    The encrypted, authenticated channel between two Pilot
    nodes. Once established, all application traffic flows
    inside it.
    \item[\textbf{Stream dial.}]
    Opening a new logical stream on top of an existing tunnel
    (analogous to opening a TCP connection to a specific port
    on an already-reachable host).
    \item[\textbf{\texttt{ErrDialToSelf}.}]
    A named error returned when a daemon tries to dial its own
    Node ID. Mirrors the same guard in \texttt{go-libp2p-swarm}.
    \item[\textbf{\texttt{viaRelay}.}]
    A per-peer boolean Pilot tracks to remember whether the
    latest authenticated packet from that peer arrived directly
    or via the beacon relay.
    \item[\textbf{Stream multiplexing (yamux).}]
    A technique for running many logical streams over one
    long-lived transport connection, with each stream opening
    in zero round-trips. Entmoot uses
    \texttt{hashicorp/yamux} to multiplex gossip frames over a
    single Pilot tunnel per peer, avoiding a per-frame stream
    dial. The same library is used by Consul, Nomad, Vault,
    and libp2p's yamux transport.
\end{description}

\subsection*{macOS operations}

\begin{description}[leftmargin=1em,style=nextline]
    \item[\textbf{Ad-hoc code signing.}]
    Signing a macOS binary with an empty identity (\texttt{-}).
    The signature is not notarised and carries no chain of
    trust, but it satisfies the local kernel's check. Go's
    cross-compiled binaries ship this way by default.
    \item[\textbf{Gatekeeper.}]
    The macOS subsystem that decides whether a binary may run.
    On Tahoe (26.2) an ad-hoc-signed binary downloaded from the
    internet triggers a stall at \texttt{\_dyld\_start} until
    Gatekeeper finishes verifying it (which, for a local daemon
    invocation, never terminates visibly).
    \item[\textbf{\texttt{xattr -cr}.}]
    Recursively removes the \texttt{com.apple.quarantine}
    extended attribute from a directory, which is what Gatekeeper
    keys its verification on.
\end{description}

\section{The Pilot fixes, in order}
\label{sec:pilot}

Pilot's patches came in three clusters:
\texttt{v1.7.2-jf.1} through \texttt{v1.7.2-jf.3} tightened the
existing transport. \texttt{v1.8.0-jf.1} and \texttt{v1.9.0-jf.1}
introduced two new pieces: a pluggable transport layer with a TCP
fallback, and a gossip-based peer-discovery overlay.
\texttt{v1.9.0-jf.2} through \texttt{v1.9.0-jf.6} are a tight
cluster of reliability fixes driven by the three-machine demo.

\subsection{\texttt{v1.7.2-jf.1}: don't punch at private addresses}

\paragraph{What broke.} When the laptop and VPS tried to
establish a tunnel, a log line appeared on the laptop:
\begin{lstlisting}
STUN returned private/unusable IP, discarding
  stun_addr=10.128.0.12
\end{lstlisting}
Shortly after, the tunnel handshake timed out. Packet captures
showed that both daemons were sending NAT-punch packets toward
\texttt{10.128.0.12}, which is a private address (see RFC 1918)
that does not exist on the public internet.

\paragraph{What it actually was.} Pilot's public STUN endpoint
sits behind a Google Cloud load balancer that rewrites every
client's source address to a private internal one
(\texttt{10.128.0.12}). The daemon had a filter that discarded
that private address when choosing its own advertised endpoint,
but \emph{the same filter was not applied to the target of
outgoing NAT-punch packets}. Every daemon was loyally sending
punch traffic to \texttt{10.128.0.12}, which is nobody's
computer.

\paragraph{How we fixed it.} \texttt{handlePunchCommand} in
\texttt{pkg/daemon/tunnel.go} now drops punches whose target
address is in an RFC 1918, loopback, link-local, or unspecified
range, logging \texttt{skipping NAT punch to non-routable target}
and falling through to the existing relay path.

\paragraph{Why the fix is right.} The condition (``is this a
routable address?'') already existed in the registration path.
The patch reuses \texttt{isPrivateAddr} symmetrically on the punch
path. Two filters with the same definition at two ingress points
is a textbook consistency fix.

\paragraph*{Tag.} \texttt{v1.7.2-jf.1}, commit \texttt{2d4e657}.
Files: \texttt{pkg/daemon/tunnel.go}.

\subsection{\texttt{v1.7.2-jf.2}: keep the tunnel alive across
rekeying}

\paragraph{What broke.} Encrypted tunnels dropped every five
minutes on the dot. Log lines like \texttt{peer marked for relay}
followed the drop; a new tunnel came back up about thirty
seconds later.

\paragraph{What it actually was.} Pilot rekeys each tunnel
periodically (part of the NetworkSync cycle). During the key
swap, the brief window where neither side knows which key to use
caused a few application-level ACKs to be dropped. The daemon's
\emph{dead-peer detector} counted those missed ACKs as silence
and concluded the peer was unreachable, tearing the tunnel down.

\paragraph{How we fixed it.} \texttt{TunnelManager} now calls a
new \texttt{notifyRekey} callback after a successful key swap.
The callback clears the ``unacknowledged keepalives'' counter
and refreshes \texttt{LastActivity} on every connection routed
over that peer, so ACKs dropped during the key swap no longer
contribute to the dead-peer tally.

\paragraph{Why the fix is right.} Dead-peer detection should
measure liveness across steady state, not across intentional
reconfigurations. The fix is a straightforward
acknowledgement-counter reset at the boundary. No
application-visible change; tunnels simply stop flapping.

\paragraph*{Tag.} \texttt{v1.7.2-jf.2}. Files: \texttt{ports.go},
\texttt{tunnel.go}, \texttt{daemon.go}.

\subsection{\texttt{v1.7.2-jf.3}: independent beacon keepalive, and
persistent visibility}

\paragraph{What broke.} Two issues. First, on consumer NATs
whose UDP idle timeout is 30 to 60 seconds, the Pilot registry
heartbeat (once every 60 seconds) was sometimes too infrequent
to keep the mapping to the beacon alive: relay traffic stopped
flowing between registry ticks. Second, running
\texttt{pilotctl set-public} (which tells the registry ``yes,
other peers may see my address'') did not survive a daemon
restart.

\paragraph{What it actually was.} The heartbeat loop was doing
double-duty and assumed both the registry and the beacon had
similar timeout budgets; they do not. Visibility state lived
only in memory, so the daemon forgot it after a restart.

\paragraph{How we fixed it.} Two additions. A new
\texttt{beaconKeepaliveLoop} runs on an independent 25-second
tick so the beacon mapping refreshes well under the strictest
consumer NAT timeout. \texttt{set-public} and the
\texttt{-public} launch flag now persist to
\texttt{\~{}/.pilot/config.json} via a new \texttt{writeConfigKey}
helper, and the daemon reads the persisted value at startup.

\paragraph{Why the fix is right.} Two different upstream services
should have two different heartbeat cadences calibrated to their
respective timeouts. Configuration that the operator set should
survive a restart. These are standard daemon-design hygiene.

\paragraph*{Tag.} \texttt{v1.7.2-jf.3}. Files: new loop in
\texttt{daemon.go}; \texttt{cmd/pilotctl/main.go}; also the
installer redirected to fetch the fork's binaries.

\subsection{\texttt{v1.8.0-jf.1}: pluggable transports and TCP
fallback}

\paragraph{What broke.} The UAE carrier was observed dropping
outbound UDP packets to the VPS's port. Pilot's entire wire
protocol was UDP. No UDP meant no tunnel, full stop.

\paragraph{What it actually was.} The daemon owned a single
\texttt{*net.UDPConn} as the only way to move bytes. There was
no mechanism to attempt an alternative protocol if UDP was
blocked.

\paragraph{How we fixed it.} Introduced a \texttt{Transport}
interface in a new \texttt{pkg/daemon/transport} package. Added
\texttt{UDPTransport} (the existing behaviour, refactored) and
\texttt{TCPTransport} (new: 4-byte big-endian length-prefix
framing, connection pooling per remote address). The registry
protocol was extended to carry a list of \texttt{endpoints} per
peer, one per transport. After direct UDP retries exhaust,
\texttt{DialConnection} attempts a TCP dial bounded by a
five-second timeout before falling through to the existing
relay path.

\paragraph{Why the fix is right.} Treating the byte-movement
layer as an interface is the canonical pattern for multi-network
systems (libp2p, Tailscale, QUIC clients). The design preserves
wire compatibility: daemons running only UDP still exchange
byte-identical frames with pre-fork versions.

\paragraph*{Tag.} \texttt{v1.8.0-jf.1}. Files: new
\texttt{pkg/daemon/transport/}, refactored \texttt{tunnel.go} and
\texttt{daemon.go}, extended \texttt{pkg/registry}.

\subsection{\texttt{v1.9.0-jf.1}: gossip-based peer discovery}

\paragraph{What broke.} The multi-transport endpoint field
needed to propagate through the registry for peers to learn
about each others' TCP listeners. The upstream public registry
does not know about the fork's new field, so the TCP fallback
was useless against an unmodified registry.

\paragraph{What it actually was.} The reachability half of peer
information (what addresses you can reach me at) was coupled to
the identity half (this node ID belongs to this key). Upgrading
reachability required an upstream cooperation we could not rely
on.

\paragraph{How we fixed it.} Added \texttt{pkg/daemon/gossip},
which exchanges \emph{signed} endpoint advertisements
peer-to-peer over existing encrypted tunnels at a 25-second
tick. The PILA handshake frame was extended with a trailing
capability bitmap, so upgraded daemons advertise gossip
capability without breaking the wire format for unmodified
upstream daemons (older peers truncate the extra bytes and
never see them).

\paragraph{Why the fix is right.} Two sources of truth (identity
anchored in the registry, reachability refreshed peer-to-peer)
is the same split libp2p's identify protocol and Tailscale's
coordination server use. Each has the right incentive: the
registry is a stable root, peer gossip is a fast-changing
directory.

\paragraph*{Tag.} \texttt{v1.9.0-jf.1}. Files: new
\texttt{pkg/daemon/gossip/}, \texttt{daemon.go} lifecycle wiring,
\texttt{pkg/registry} snapshot format.

\subsection{\texttt{v1.9.0-jf.2}: NAT-drift refresh and macOS
code-sign workaround}

\paragraph{What broke.} Two seemingly unrelated symptoms. On the
laptop, outbound stream dials over an otherwise-healthy tunnel
timed out with ``dial timeout'' even though \texttt{pilotctl
peers} reported the peer as encrypted and authenticated. On the
same laptop, interactively-launched \texttt{pilotctl info} hung
indefinitely at \texttt{\_dyld\_start}.

\paragraph{What it actually was.} The laptop's UAE carrier was
rotating its external source port between rekeys. Pilot cached
the peer's UDP address on key exchange and never refreshed it
afterwards, so the VPS kept sending to the stale port. The
carrier's NAT silently dropped everything to that old port.
Meanwhile, macOS 26.2 Tahoe's Gatekeeper was stalling every
interactively-launched ad-hoc-signed Go binary in early dynamic
loader initialisation.

\paragraph{How we fixed it.} Two changes. The daemon refreshes
its cached peer endpoint inside \texttt{handleEncrypted} on any
successful authenticated decrypt (guarded by a read-lock
pre-check so the write path is not churned). The installer now
re-signs binaries with the local ad-hoc identity and clears
\texttt{com.apple.quarantine} extended attributes on macOS,
which satisfies Gatekeeper's first-launch gate.

\paragraph{Why the fix is right.} The cached endpoint is the
daemon's latest evidence of reachability; every authenticated
packet is fresh evidence. The endpoint should track the evidence.
The macOS workaround is what every cross-compiled Go tool needs
under Tahoe; adding it to the installer once beats every
operator discovering it the hard way.

\paragraph*{Tag.} \texttt{v1.9.0-jf.2}. Files: \texttt{tunnel.go},
\texttt{install.sh}.

\subsection{\texttt{v1.9.0-jf.3}: automatic recovery from
one-sided restart}

\paragraph{What broke.} When the VPS daemon restarted, its
in-memory session keys were gone, but the laptop still had the
old keys and kept sending encrypted traffic. The VPS logged
\begin{lstlisting}
encrypted packet from node but no key peer_node_id=45491
\end{lstlisting}
forever. The two sides never spontaneously re-handshook. A
manual daemon restart on the laptop was the only way out.

\paragraph{What it actually was.} The receiving side could tell
it had no key for a peer, but it had no policy to do anything
about it. The sending side saw its own keepalive ACKs succeeding
at the tunnel level and had no way to observe the other side's
confusion.

\paragraph{How we fixed it.} On any decrypt failure caused by
missing crypto state, the daemon emits an unsolicited PILA
(handshake initiation) to the source address of the offending
frame, rate-limited to once per minute per peer (so a hostile
sender cannot amplify the mechanism into a reflection attack).
Recovery completes within 200 to 300 milliseconds of the peer's
next keepalive.

\paragraph{Why the fix is right.} The fix puts the observation
and the remediation on the same side: the node that can detect
the desync is the one that asks for a re-handshake. The
rate-limit prevents abuse. Equivalent to the ``force rekey on
decrypt fail'' rule WireGuard applies to its roaming handshake.

\paragraph*{Tag.} \texttt{v1.9.0-jf.3}. Files: \texttt{tunnel.go}
(\texttt{maybeSendRecoveryPILA}).

\subsection{\texttt{v1.9.0-jf.4}: reply on the path the request
came in on}

\paragraph{What broke.} Tunnels showed
\texttt{encrypted=yes auth=yes relay=true} in the peer table, and
keepalives flowed both ways, but any attempt to open a new
application-layer stream timed out. This affected every
VPS-to-laptop dial for the whole day.

\paragraph{What it actually was.} Pilot kept two maps: one of
direct UDP addresses (\texttt{peers}) and a boolean flag
(\texttt{relayPeers}) saying whether a peer was reachable via
relay. The two could disagree across daemons: A's flag could
say ``via relay'' while B's map still held a stale direct
address for A. Worse, when a relay frame arrived from an
unknown peer, the daemon cached the beacon's address as the
peer's direct address, which then polluted same-LAN detection,
the freshly-added NAT-drift refresh, and the direct-UDP fallback
path.

\paragraph{How we fixed it.} Replaced both maps with a single
\texttt{paths map[uint32]*peerPath}. Each entry records the
\emph{most recent authenticated ingress path}: either a direct
UDP address or ``via relay''. Every authenticated decrypt
updates the entry. Every outbound reply is routed on the same
path the last request arrived on.

Also: relay round-trip time is two to three times direct, so the
SYN retransmission budget starts at three seconds when the peer
is in relay mode (up from one second).

\paragraph{Why the fix is right.} This is the ``reply on
ingress'' pattern WireGuard uses for roaming, and what Tailscale
calls ``DERP sticky''. Symmetry between the two daemons falls out
of each one observing its own ingress traffic, instead of
requiring both sides to agree on a shared model of the network.

\paragraph*{Tag.} \texttt{v1.9.0-jf.4}. Files: \texttt{tunnel.go},
\texttt{tunnel\_path\_test.go}.

\subsection{\texttt{v1.9.0-jf.5}: clean up all relay state}

\paragraph{What broke.} A VPS restart on
\texttt{v1.9.0-jf.4} left relay-only peers in the
``encrypted packet from no key'' loop indefinitely. The jf.3
auto-recovery fired, but the unsolicited PILA went out to a
zero-address placeholder and was dropped silently.

\paragraph{What it actually was.} Several per-peer state
structures were not fully cleared across \texttt{AddPeer} or
\texttt{RemovePeer}, so stale values from prior sessions
poisoned the new session. Also, \texttt{maybeSendRecoveryPILA}
unconditionally wrote to the direct UDP address, even when the
triggering frame had arrived via relay.

\paragraph{How we fixed it.} Five small changes, all about
per-peer hygiene:
\begin{itemize}[leftmargin=1.2em]
    \item Recovery PILA now routes through the beacon when the
    triggering frame arrived via relay.
    \item \texttt{AddPeer} clears \texttt{viaRelay} so a fresh
    direct endpoint is authoritative.
    \item \texttt{RemovePeer} wipes all per-peer state, not
    just two of the seven structures.
    \item \texttt{handlePacket} no longer auto-adds a peer under
    the relay-origin zero-address placeholder.
    \item \texttt{maybeSendRecoveryPILA} checks the UDP socket
    is non-nil before writing (defensive guard during startup).
\end{itemize}

\paragraph{Why the fix is right.} \texttt{RemovePeer} should be
``reset all per-peer state'', not ``reset some of it''. This is
WireGuard's \texttt{wg set peer \ldots remove} contract. The
PILA routing fix matches the same reply-on-ingress principle
introduced in jf.4, applied to the recovery path.

\paragraph*{Tag.} \texttt{v1.9.0-jf.5}. Files: \texttt{tunnel.go},
\texttt{daemon.go}.

\subsection{\texttt{v1.9.0-jf.6}: don't dial yourself}

\paragraph{What broke.} Gossip propagation inflated from
sub-second to minute-scale latencies. The VPS's
\texttt{pilot-daemon} was burning 210\% CPU (two cores) and
pushing roughly 5,900 packets per second, most of them looping
to itself.

\paragraph{What it actually was.} Every Entmoot invite minted to
that point listed the issuer as a bootstrap peer, including on
the issuer's own machine. Pilot's \texttt{ensureTunnel} (callable
from code paths other than the Entmoot bootstrap filter) had no
guard against dialling the local Node ID. On a multi-homed host
the ``same-LAN peer detected'' logic then matched both a Docker
bridge address (\texttt{172.17.0.1}) and the public IP as valid
local endpoints, so the daemon established multiple duplicate
self-tunnels, each happily retransmitting into the other.

\paragraph{How we fixed it.} \texttt{ensureTunnel} and
\texttt{DialConnection} now short-circuit with a typed sentinel
\texttt{protocol.ErrDialToSelf} when given the local Node ID.
The invite minter was also fixed (see Entmoot \texttt{v1.0.8}).

\paragraph{Why the fix is right.} \texttt{ErrDialToSelf} is the
canonical name for this guard (go-libp2p swarm, for one, uses
the same pattern). Returning a typed sentinel rather than
silently succeeding makes caller-side invariant violations
visible. Defence in depth: the fix lands at both the caller (the
invite minter) and the callee (the daemon), so a future invite
that happens to include self still cannot cause a storm.

\paragraph*{Tag.} \texttt{v1.9.0-jf.6}. Files: \texttt{daemon.go}.

\section{The Entmoot fixes, in order}
\label{sec:entmoot}

Entmoot's releases are easier to summarise: each one was driven
by a symptom on the live demo, and the fixes moved steadily from
``make push-only gossip reliable'' (v1.0.2, v1.0.3) to
``implement the canonical gossip design'' (v1.0.4, v1.0.5) to
``tune it for real networks'' (v1.0.6, v1.0.7, v1.0.8) to
``remove per-frame transport overhead'' (v1.1.0).

\subsection{\texttt{v1.0.0}: the foundations}

\paragraph*{Not strictly a fix;} this is the shape of v1 that
the later releases patch. Five points to remember:

\begin{itemize}[leftmargin=1.2em]
    \item \textbf{SQLite store.} One database per group,
    write-ahead logging, one database file per group's
    messages. Replaces the JSONL backend used in pre-v1
    development.
    \item \textbf{Control-socket IPC.} A single long-running
    \texttt{entmootd join} process per host owns the Pilot
    connection and is the only writer to the store. Short-lived
    CLI invocations (\texttt{publish}, \texttt{tail},
    \texttt{info}) reach it through a Unix-domain socket at
    \texttt{\~{}/.entmoot/control.sock}.
    \item \textbf{Five-command CLI.} \texttt{join},
    \texttt{publish}, \texttt{tail}, \texttt{query},
    \texttt{info}. \texttt{join} is the only blocking one.
    \item \textbf{Invites with TTL.} A new \texttt{ValidUntil}
    field in the invite bundle; default 24 hours. Expired
    invites are rejected.
    \item \textbf{Installer and skill.}
    \texttt{install.sh} (prebuilt binary fast path, source fall
    back) and a \texttt{SKILL.md} for agent-driven installation.
\end{itemize}

\paragraph*{Tag.} \texttt{v1.0.0}, 17 April 2026.

\subsection{\texttt{v1.0.2}: retries and anti-entropy}

\paragraph{What broke.} Push-only gossip from v1.0.0 dropped any
message it failed to deliver on the first try. A peer that was
briefly offline silently missed messages forever; the only
recovery was a full re-join.

\paragraph{What it actually was.} The design had no concept of
``try again later'' and no background synchronisation to paper
over gaps that the push layer missed.

\paragraph{How we fixed it.} Two new mechanisms:

\begin{itemize}[leftmargin=1.2em]
    \item An exponential-backoff retry loop. Transient push and
    fetch failures are queued; a \texttt{retryLoop} goroutine
    drains the pending queue every 500 ms with delays
    \texttt{1s, 2s, 4s, 8s, 16s, 30s, 60s \(\times\) 4}.
    Exhausted slots are logged and left to the next mechanism.
    \item Anti-entropy reconciliation. When a peer reconnects
    (either as inbound accept or as a successful outbound dial),
    the side that detects the reconnection starts a round with
    it: request the peer's Merkle root, compare with local, and
    if they differ, request the list of message IDs since the
    local latest timestamp and fetch the missing bodies.
    Rate-limited to one round per peer per minute.
\end{itemize}

Two new wire types for the range enumeration:
\texttt{MsgRangeReq} (\texttt{0x09}) and \texttt{MsgRangeResp}
(\texttt{0x0A}).

\paragraph{Why the fix is right.} Every production gossip
system has both a fast path (push with retry) and a slow path
(anti-entropy). Cassandra, Riak, Scuttlebutt: all pair the two.
Push handles the common case; anti-entropy is the safety net.

\paragraph*{Tag.} \texttt{v1.0.2}, 19 April 2026.

\subsection{\texttt{v1.0.3}: publish returns immediately}

\paragraph{What broke.} \texttt{entmootd publish} took 30
seconds to return on every call, and the return always printed
an IPC read-response timeout, even when the message had been
stored locally and was propagating normally.

\paragraph{What it actually was.} The publish handler tried to
deliver the message to all peers synchronously before
acknowledging the CLI. One slow peer blocked the whole command.

\paragraph{How we fixed it.} The handler now returns as soon as
the message is durably written to the local store. Peer
delivery happens asynchronously through the retry scheduler.
Publish latency drops to single-digit milliseconds regardless of
peer reachability.

\paragraph{Why the fix is right.} This is the industry-standard
publish contract. Kafka durable-log commit, NATS fire-and-forget,
RabbitMQ publisher confirms: ``publish'' means ``accepted for
delivery'', not ``delivered to every subscriber''. The retry
scheduler and anti-entropy reconciliation from v1.0.2 already
carry the delivery guarantee.

\paragraph*{Tag.} \texttt{v1.0.3}, 19 April 2026.

\subsection{\texttt{v1.0.4}: Plumtree}

\paragraph{What broke.} Push-only gossip struggled under
asymmetric reachability. On the canary, the laptop and Phobos
had no direct edge: neither could reach the other at the
transport layer. A message from Phobos reached the VPS but
never reached the laptop, because Entmoot had no re-forwarding
step.

\paragraph{What it actually was.} The gossip design was
``publisher fans out; receivers store and stop''. A partial
topology guaranteed gaps that only anti-entropy could eventually
repair, at the cost of up to a minute of delay.

\paragraph{How we fixed it.} Switched to a Plumtree-style
protocol (Leit\~ao, Pereira, Rodrigues, SRDS 2007).

Each peer sits in exactly one of two sets per message:
\texttt{eagerPushPeers} (get the full body up front) or
\texttt{lazyPushPeers} (get only the ID). Three new control
messages maintain this split:

\begin{itemize}[leftmargin=1.2em]
    \item \texttt{IHAVE} (\texttt{0x0B}): ``I have message X, but
    I am sending only the ID''.
    \item \texttt{GRAFT} (\texttt{0x0C}): ``I heard IHAVE for X,
    did not receive the body within three seconds, please send
    it and promote me to your eager set''.
    \item \texttt{PRUNE} (\texttt{0x0D}): ``I already had X via
    another path; drop me from your eager set to avoid
    duplicates''.
\end{itemize}

Initial state: everyone is eager. Duplicates self-prune into a
spanning subset per message. On the canary, a Phobos publish now
reaches the laptop via VPS as the re-fanout hop.

\paragraph{Why the fix is right.} Plumtree is the canonical way
to build a self-healing broadcast tree. Its overhead at steady
state is one push per spanning-tree edge, instead of the
multiplicative explosion of naive gossip. Riak Core, Helium,
Solana, libp2p GossipSub all use variants.

\paragraph*{Tag.} \texttt{v1.0.4}, 20 April 2026.

\subsection{\texttt{v1.0.5}: re-fanout on every path}

\paragraph{What broke.} After v1.0.4, re-fanout fired only when
a message arrived via the live push path. Messages acquired via
retry or anti-entropy reconciliation were stored but not
forwarded.

\paragraph{What it actually was.} The ``first-seen forwarding''
rule (forward exactly the first time, regardless of source) was
implemented in only one of three acquisition call sites. The
other two silently broke it.

\paragraph{How we fixed it.} Centralised the
\texttt{plumCancelGraftsFor(id) + refanout(peer, id)} hook
inside \texttt{fetchFrom} after a successful \texttt{Store.Put}.
All three call sites inherit it.

\paragraph{Why the fix is right.} The rule libp2p GossipSub
applies: ``not seen before then forward to mesh, irrespective of
arrival method''. One place in the code, one behaviour.

\paragraph*{Tag.} \texttt{v1.0.5}, 20 April 2026.

\subsection{\texttt{v1.0.6}: trust-aware fanout, parallel sends,
dial-backoff}

\paragraph{What broke.} Cross-mesh publish latency regressed to
about four minutes. Fanout was iterating over every peer
serially; each dial to an unreachable peer paid Pilot's full
32-second timeout; and the retry scheduler dialled the same
dead peer again for every pending message.

\paragraph{What it actually was.} Three independent gaps
compounded. Fanout did not ask ``is this peer even reachable''
before attempting. Fanout was sequential so one slow peer
blocked everyone. There was no backoff memory so the same dead
peer was dialled over and over.

\paragraph{How we fixed it.} Three coordinated changes:

\begin{itemize}[leftmargin=1.2em]
    \item \textbf{Trust-aware reachability oracle.} Before
    building the fanout set, call Pilot's
    \texttt{TrustedPeers} and filter. Peers with no trust pair
    (structurally unreachable for the life of this config) are
    dropped up-front. Cached for one second so bursty fanouts
    do not ask Pilot per message per peer.
    \item \textbf{Parallel fanout with per-peer timeout.} One
    goroutine per peer, each with a five-second
    \texttt{context.WithTimeout} and a concurrency cap of 32.
    One slow peer no longer blocks the others.
    \item \textbf{Per-peer dial-backoff cache.} After one
    failed dial, further dials short-circuit for an
    exponentially expanding window (one second, two, four,
    eight, capped at five minutes). A successful dial clears the
    window.
\end{itemize}

Together they collapse cross-mesh latency from roughly four
minutes back to sub-second.

\paragraph{Why the fix is right.} Each of the three changes is
the canonical pattern from a mature system: the reachability
oracle is how libp2p GossipSub treats connectedness; parallel
fanout is idiomatic Go (\texttt{errgroup} with
\texttt{SetLimit}); the dial-backoff cache is how
\texttt{go-libp2p swarm.DialBackoff} works. Together they are
the standard triptych for making ``don't spend real time on a
dead peer'' work.

\paragraph*{Tag.} \texttt{v1.0.6}, 20 April 2026.

\subsection{\texttt{v1.0.7}: inline bodies and jittered backoff}

\paragraph{What broke.} Even after v1.0.6, a 1.5-minute
propagation tail persisted. Each gossip hop paid two separate
Pilot stream dials: sender's push, receiver's fetch-back for
the body. And the deterministic retry schedule synchronised
every node's retries to the same wall-clock moments, so a
single NAT flap cascaded into many minutes of propagation
delay.

\paragraph{What it actually was.} The push layer had shipped
only the message ID, expecting the receiver to pull the body.
The receiver's pull was a whole new stream dial. And exact
retry delays made independent nodes retry in lockstep.

\paragraph{How we fixed it.} Two changes:

\begin{itemize}[leftmargin=1.2em]
    \item \textbf{Inline body on eager push.} Gossip frames for
    messages smaller than 4 KiB now carry the full body.
    Receivers hash-verify and store directly, skipping the
    fetch-back round-trip. This is what the canonical Plumtree
    paper specified in the first place; every surveyed
    production gossip (libp2p GossipSub, Scuttlebutt EBT,
    Matrix federation, Cassandra, IPFS Bitswap) does this.
    \item \textbf{Decorrelated-jitter retry backoff.} Replaces
    the fixed schedule with \texttt{next = min(cap,
    random(base, prev $\times$ 3))}, Marc Brooker's 2015 AWS
    formula. Same 10-attempt budget, but nodes no longer retry
    in synchronised waves.
\end{itemize}

The reconcile cooldown also dropped from 60 seconds to a
jittered 30 seconds (that range matches what GossipSub and
Plumtree \S3.4 use for lazy-push cadence, and is an order of
magnitude more appropriate than what had been in place).

\paragraph{Why the fix is right.} Inline bodies bring Entmoot
back to what the Plumtree paper actually described; Entmoot's
earlier ID-only fanout was a non-canonical simplification.
Decorrelated jitter is the specific, named cure for
retry-synchronisation cascades, which is exactly what we were
seeing.

\paragraph*{Tag.} \texttt{v1.0.7}, 21 April 2026.

\subsection{\texttt{v1.0.8}: invites should not include the
issuer}

\paragraph{What broke.} See the Pilot
\texttt{v1.9.0-jf.6} entry above: self-amplified packet storm
on the VPS, 5,900 packets per second, 210\% CPU.

\paragraph{What it actually was.} Every invite minted up to
that point listed the issuer as a bootstrap peer. When the
issuer received its own invite through ambient membership
channels, the Pilot layer happily established self-tunnels.

\paragraph{How we fixed it.} Two defence-in-depth changes:

\begin{itemize}[leftmargin=1.2em]
    \item \texttt{entmootd invite create} excludes the issuer's
    Node ID from \texttt{bootstrap\_peers} at mint time.
    \item Pilot \texttt{v1.9.0-jf.6} adds the
    \texttt{ErrDialToSelf} sentinel so that any future invite
    that happens to include self cannot cause a storm.
\end{itemize}

Legacy invites still work: the \texttt{Gossiper.Join} path
already has a self-skip filter at every bootstrap strategy. The
new defences close the remaining holes.

\paragraph{Why the fix is right.} Cassandra's gossiper
maintains the invariant that \texttt{live\_endpoints} excludes
self by construction. This is the same invariant, applied at
mint time. Defence in depth: fix the source, and also fix the
callee so a future regression cannot amplify.

\paragraph*{Tag.} \texttt{v1.0.8}, 21 April 2026.

\subsection{\texttt{v1.1.0}: one persistent session per peer}

\paragraph{What broke.} Both peers on a node were observed
failing simultaneously at the Pilot stream-dial layer, even while
the underlying tunnel reported \texttt{encrypted=yes auth=yes}.
Each dropped dial cost Pilot's full roughly 32-second timeout
before gossip could give up and try a retry; with two peers
wedged in parallel, the accumulated timeouts dominated
propagation latency.

\paragraph{What it actually was.} Every gossip frame was opening
a fresh Pilot stream, paying the full SYN handshake per frame.
On a healthy tunnel that overhead is negligible, but whenever a
stream SYN/ACK silently wedged (middlebox, NAT timer expiry,
any of several causes we had chased in the Pilot layer), gossip
was exposed to the 32-second Pilot dial budget on the new
stream, even though the tunnel itself was fine.

\paragraph{How we fixed it.} The Pilot transport adapter now
maintains one long-lived \texttt{hashicorp/yamux} session per
peer and multiplexes every gossip stream over it. The first
contact still dials a Pilot stream; subsequent frames open a
yamux stream on the cached session with no extra handshake
(zero round-trip new-stream). yamux's built-in keepalive (30-second
interval, 10-second timeout) detects silently-wedged sessions
and closes them; the next dial opens a fresh session
automatically.

Wire compatibility: Entmoot's JSON frame format is unchanged;
yamux simply multiplexes the same frames over the session.
Mixed v1.0.x and v1.1.0 peers do \emph{not} interoperate
(v1.0.x expects a raw stream per frame; v1.1.0 expects a yamux
session), so every node in a group must upgrade together.

\paragraph{Why the fix is right.} Multiplexed long-lived
sessions are how every mature stream-based overlay handles the
per-frame overhead problem: Consul, Nomad, Vault, and libp2p's
yamux transport all use the same library. Pushing the keepalive
down into the multiplexer means the gossip layer does not have
to build its own liveness check for sessions, and silently-wedged
transports are detected and recovered without any gossip-layer
plumbing.

\paragraph*{Tag.} \texttt{v1.1.0}, 21 April 2026. Files:
\texttt{src/pkg/entmoot/transport/pilot/pilot.go},
new test \texttt{pilot\_yamux\_test.go}, new dependency
\texttt{github.com/hashicorp/yamux}.

\subsection{\texttt{v1.2.0}: transport-endpoint advertisements
over Entmoot gossip}

\paragraph{What broke.} A three-peer live test designed to
exercise Pilot's TCP-fallback transport stack (which would let
Entmoot keep running when UDP is blocked) showed that the
fallback code path was operationally inert --- the Pilot daemon's
multi-transport dialer had a valid TCP path compiled in, but
\texttt{peerTCP} was always empty at dial time, so the dialer
always fell through to the beacon-relay path instead.

\paragraph{What it actually was.} Our fork's daemon dutifully
advertises its TCP endpoint to the central registry at register
time. The upstream stock registry at \texttt{34.71.57.205:9000}
silently drops the \texttt{endpoints} field in the register
frame, so peers' registry lookups return only the UDP
\texttt{real\_addr}. The field-drop is an upstream bug we cannot
patch; any attempt to route around it by standing up a
fork-built registry would force a mesh-wide rebuild every time
someone on the fork pushed a new tag.

\paragraph{How we fixed it.} We distributed per-peer TCP
endpoints \emph{through Entmoot itself}. Each member signs a
small \texttt{TransportAd} record naming its current endpoints
and gossips it through the existing Plumtree fanout. Receivers
verify the signature against the roster entry, store-or-replace
under Last-Writer-Wins semantics keyed on
\texttt{(group, author)}, and install the endpoints into Pilot's
\texttt{peerTCP} cache via a new \texttt{SetPeerEndpoints} IPC
command (Pilot \texttt{v1.9.0-jf.7}). The record is bounded at
one row per \texttt{(group, member)} regardless of publish rate,
and three layers of spam defence run in cost order on receive:
schema-and-size validation ($\le 1$~KiB), publisher allowlist
(sender \emph{must} equal author), and per-(peer, topic)
token-bucket rate limit. Only after all three pass does the
Ed25519 signature verification run. A latent second bug in
anti-entropy also surfaced during the test but was deferred to
\texttt{v1.2.1}: \texttt{fetchPeerRange} used
\texttt{since = latestLocalTimestamp} as its lower bound, so
gaps older than the local node's newest message were silently
excluded from the range query.

\paragraph{Why the fix is right.} Routing control-plane state
through the data-plane gossip channel is a canonical 2020s
federated-protocol move: Matrix uses signed state events over
the same federation stream that carries message content;
IPNS-over-pubsub distributes name pointers over libp2p pubsub
rather than a separate DHT write path. The LWW +
publisher-allowlist pattern is the structural spam bound that
Matrix room ACLs and GossipSub v1.1 peer scoring both converge
on. Retaining one row per (group, author) keeps storage
structurally bounded at $O(N)$ in the member count regardless
of publish rate --- the same invariant IPNS records enforce
for distributed name pointers.

\paragraph*{Tag.} \texttt{v1.2.0}, 21 April 2026. Files:
\texttt{src/pkg/entmoot/wire/types.go},
\texttt{src/pkg/entmoot/store/sqlite.go} (new
\texttt{transport\_ads} table),
\texttt{src/pkg/entmoot/gossip/gossiper.go} (new
\texttt{onTransportAd} handler),
\texttt{src/pkg/entmoot/transport/pilot/pilot.go} (new
\texttt{SetPeerEndpoints} call), Pilot
\texttt{pkg/driver/driver.go} + \texttt{pkg/daemon/ipc.go}.

\subsection{\texttt{v1.2.1}: Range-Based Set Reconciliation,
background anti-entropy, and tunnel-up trigger}

\paragraph{What broke.} A second live test surfaced the
anti-entropy bug that \texttt{v1.2.0} had merely postponed. Two
messages published from one peer (Phobos, Italian residential
NAT) during a brief reachability outage never propagated to the
public VPS, even after the VPS's reconcile fired on an unrelated
fresh message from the same peer. The \texttt{MerkleReq}
short-circuit correctly reported roots mismatched; the
subsequent fetch step still couldn't bridge the gap.

\paragraph{What it actually was.} Two orthogonal causes working
together. The first: \texttt{fetchPeerRange(since = t)} used
the local node's newest-message timestamp as the range's lower
bound, so any gap with a timestamp \emph{older} than the local
newest was excluded from the query. The second: reconciliation
fired only on inbound gossip events, not on a background timer,
so a quiet partition with no organic publish traffic could leave
a gap sitting indefinitely even after both sides reconnected.

\paragraph{How we fixed it.} Three changes landed together. The
reconcile algorithm moved from the timestamp-cursor fetch to
Range-Based Set Reconciliation (RBSR). RBSR recursively splits
the 32-byte message-identifier keyspace into sub-ranges,
exchanges a 16-byte fingerprint per sub-range, and descends only
where fingerprints disagree. Bandwidth is proportional to the
symmetric difference, convergence is $O(\log N)$ rounds, and
gaps anywhere in the timeline are found by construction. The
Pilot-transport adapter gained a \texttt{SetOnTunnelUp(cb)}
callback that fires \texttt{maybeReconcile} within milliseconds
of every newly-established yamux session (outbound and inbound
both). And a background \texttt{reconcilerLoop} goroutine ticks
every 30\,s $\pm\,20\%$ jitter, picks the least-recently-
reconciled roster member, and silently skips when the peer's
cached Merkle root still matches local --- an idle mesh pays
zero dials per tick at steady state.

\paragraph{Why the fix is right.} RBSR is the canonical 2020s
set-diff algorithm; Willow, Earthstar, and iroh all adopt it,
as does the Nostr Negentropy protocol (NIP-77). The 30\,s
background ticker
matches Serf's \texttt{PushPullInterval} default for LAN
gossip, chosen for exactly this ``partition healed but no
organic publish'' scenario. Tunnel-up event triggers match
Secure Scuttlebutt's Epidemic Broadcast Trees
on-connect handshake and Automerge sync's
have-heads-on-connect pattern; tying anti-entropy to
connection events, not just timers, lets a flaky network
converge as fast as its reconnects allow. Wire-compatibility
with \texttt{v1.2.0} peers is preserved: the legacy
\texttt{MsgRangeReq} and \texttt{MsgRangeResp} handlers stay
active, and a \texttt{v1.2.1} initiator that encounters a
\texttt{v1.2.0} responder observes an \texttt{EOF} and falls
back silently.

\paragraph*{Live field verification (2026-04-22).} Re-running
the original bug pattern on the live three-peer mesh with all
peers on \texttt{v1.2.1}: VPS $\leftrightarrow$ laptop converged
in three RBSR rounds with \texttt{missing}~$=0$. VPS
$\leftrightarrow$ Phobos, after Phobos cold-restarted its Pilot
daemon to clear the pre-existing tunnel wedge, converged in
three rounds with \texttt{missing}~$=13$ (Phobos had
accumulated thirteen messages across partition events); the
VPS store grew from 95 to 108 messages, an exact match on the
RBSR-reported difference. The \texttt{OnTunnelUp} callback
fired 2.1\,s after the yamux session became usable. Zero
\texttt{session ended early} frames were observed, confirming
the wire-compat fallback stayed dormant once the whole mesh was
on \texttt{v1.2.1}. After convergence, the idle-mesh
\texttt{tick skip (roots match cache)} debug log fired on the
background ticker, confirming the cached-root-match
optimization kept the quiescent mesh silent.

\paragraph*{Tag.} \texttt{v1.2.1}, 22 April 2026. Files: new
package \texttt{src/pkg/entmoot/reconcile/} (the RBSR state
machine), \texttt{src/pkg/entmoot/wire/types.go} (new
\texttt{MsgReconcile} opcode),
\texttt{src/pkg/entmoot/gossip/gossiper.go} (rewritten
\texttt{reconcileWith}, new \texttt{onReconcileReq} and
\texttt{reconcilerLoop}),
\texttt{src/pkg/entmoot/gossip/transport.go} +
\texttt{src/pkg/entmoot/transport/pilot/pilot.go}
(\texttt{SetOnTunnelUp}),
\texttt{src/pkg/entmoot/store/\{store,memory,sqlite,jsonl\}.go}
(new \texttt{IterMessageIDsInIDRange}). No Pilot change.

\section{The second shakedown: hide-IP, TURN, and port 1004}
\label{sec:secondshakedown}

The first shakedown ended with normal-mode gossip converging. The
second began when we asked a harder question: can the laptop hide its
real address and still participate in the mesh? That moved the problem
from ordinary NAT traversal to strict TURN-mediated operation. In
normal Pilot mode a peer may try direct UDP, TCP, beacon relay, cached
connections, and peer TURN endpoints. In strict
\texttt{-outbound-turn-only} mode the rule is different: every
outbound peer packet must leave through the local node's own TURN
allocation, and if that path cannot send, Pilot must fail closed rather
than leak the real address.

\subsection{Pilot \texttt{jf.8}--\texttt{jf.11b}: TURN and push
rotation}

\paragraph{What broke.} The laptop could publish a TURN address, but
the rest of the mesh kept learning stale relay ports after Pilot
restarts or Cloudflare allocation rotations. Entmoot's first
implementation polled \texttt{pilotctl info} every 30 seconds and
republished when the address changed. That worked, but it was both
slow and wasteful: the common case is no rotation.

\paragraph{What it actually was.} TURN relay addresses are state, not
configuration. They are assigned by the provider and can rotate under
the daemon. Treating them like a static \texttt{-advertise-endpoint}
value left peers pinned to dead allocations.

\paragraph{How we fixed it.} Pilot \texttt{jf.8} added TURN client
support and credential providers. The \texttt{jf.11a} series added the
privacy posture flags and \texttt{Info} fields that let Entmoot verify
\texttt{-hide-ip}. Pilot \texttt{jf.11b} added IPC
\texttt{Subscribe}/\texttt{Notify} and a \texttt{turn\_endpoint} topic.
Entmoot \texttt{v1.4.0} added \texttt{-hide-ip} transport ads, and
\texttt{v1.5.0} replaced steady-state polling with a single
subscription. Old Pilot daemons still fall back to the 30-second poller.

\paragraph{Why the fix is right.} The application layer needs to
publish current reachability, but it should not have to poll a local
daemon forever. Push turns a rare state change into a rare IPC event.
It also makes the privacy check explicit: Entmoot can suppress
advertised addresses, but only Pilot can prevent outbound source-IP
leaks.

\subsection{Pilot \texttt{jf.12}--\texttt{jf.15.7}: rendezvous and
bilateral TURN}

\paragraph{What broke.} Once all three machines ran TURN-aware builds,
the next deadlock was subtler. VPS knew the laptop's new TURN
allocation, but the laptop had not installed the reciprocal permission
for the VPS allocation. Packets left one side and disappeared at the
other side's TURN permission check. Because UDP send APIs report only
local send success, both daemons could appear healthy while no bytes
arrived.

\paragraph{What it actually was.} TURN permissions are bilateral and
allocation-local. Installing a peer's endpoint once is not enough: a
local allocation rotation can erase the permission, and a peer rotation
can make a cached permission stale. The push path itself depends on the
data path, so it cannot be the only recovery mechanism.

\paragraph{How we fixed it.} Pilot \texttt{jf.14} introduced a small
Pkarr-style rendezvous service for signed TURN endpoint records;
\texttt{jf.14.2} added refresh and stale-cache eviction;
\texttt{jf.15.3} installed reciprocal \texttt{CreatePermission} state
when peer TURN endpoints are installed; \texttt{jf.15.4} made the
rendezvous lookup fire in strict TURN-only mode; \texttt{jf.15.6}
allowed recovery PILA to fire for TURN-delivered frames whose observed
remote address is not a direct UDP endpoint; and \texttt{jf.15.7}
added periodic peer re-lookup so a data-path break does not prevent
endpoint refresh. Entmoot's role in this phase was to keep publishing
current transport ads, replay persisted ads at startup, and expose
enough trace signal to separate discovery failure from transport
failure.

\paragraph{Why the fix is right.} Rendezvous is discovery, not routing.
It gives each Pilot daemon the fresh endpoint material needed to
install permissions and evict stale cached connections. The actual
route decision remains in Pilot, where \texttt{-outbound-turn-only} can
win deterministically over direct UDP, TCP, beacon relay, and peer-TURN
preferences.

\subsection{Entmoot \texttt{v1.5.1}--\texttt{v1.5.10}: stale streams
and reconcile recovery}

\paragraph{What broke.} After Pilot peers finally reached
\texttt{encrypted=yes auth=yes}, Entmoot still failed to converge.
VPS could open or attempt streams to laptop port~1004, but reconcile
sessions ended early, timed out, or retried through the same dead
yamux/Pilot stream. Store counts diverged even though the underlying
mesh was alive.

\paragraph{What it actually was.} Several bugs stacked in the same
symptom. The responder side of RBSR could learn that it was missing
message IDs but fail to fetch the bodies. Short push contexts leaked
into background reconcile attempts. Request/response deadlines started
too early, before Pilot stream establishment completed. Stale
connection IDs from the Pilot IPC layer looked like generic dial
failures, which poisoned Entmoot's backoff instead of evicting the
cached session. Finally, closing or closing-soon Pilot streams could be
accepted into yamux and reused until higher-level timeouts fired.

\paragraph{How we fixed it.} Entmoot \texttt{v1.5.1} made responder
RBSR fetch symmetric. \texttt{v1.5.2} through \texttt{v1.5.4} split
reconcile attempt lifetimes from push lifetimes, separated Pilot dial
budget from wire I/O deadlines, and evicted stale sessions on response
timeouts. \texttt{v1.5.5} preserved early Pilot IPC close frames so a
freshly registered connection observes EOF immediately. \texttt{v1.5.6}
fixed transport-ad self-install and forwarded snapshot validation.
\texttt{v1.5.8} retried stale Pilot stream IDs without poisoning
dial-backoff. \texttt{v1.5.9} introduced an explicit Pilot/yamux
session lifecycle manager with generation ids, draining, active-stream
counts, and a \texttt{-trace-gossip-transport} mode. \texttt{v1.5.10}
added \texttt{-trace-reconcile}, Pilot readiness waiting at startup,
and transport-owned stream error classification.

\paragraph{Why the fix is right.} Reconcile is a correctness path, not
a best-effort push. It must treat ``the cached stream died'' as a local
session repair event, not as proof that the peer is unreachable. The
current split is the useful one: gossip owns retry policy and
anti-entropy semantics; the Pilot adapter owns IPC/yamux-specific error
classification and session lifecycle.

\subsection{Pilot \texttt{jf.15.8}--\texttt{jf.15.21}: the stream
layer becomes explicit}

\paragraph{What broke.} The last visible symptom was the famous
\texttt{:1004} timeout. Small frames could pass, but large reconcile
responses stalled, retransmitted until max retry, and reset. Other
failures came from accepting half-open streams before the final ACK,
from stale daemon-side stream ids, or from listener ports that survived
an Entmoot restart.

\paragraph{What it actually was.} TURN routing had become mostly
correct; the remaining bugs were in Pilot's virtual stream layer. The
old implementation mixed SYN/SYN-ACK/ACK, port listener dispatch, IPC
accept, retransmit state, TIME\_WAIT cleanup, and send errors in one
large path. That made symptoms look like Entmoot reconcile bugs even
when the failed invariant was lower-level stream lifecycle.

\paragraph{How we fixed it.} Pilot added \texttt{-trace-streams},
explicit stream transition/removal logs, listener \texttt{Unbind}, a
passive-open rule that exposes accepted streams only after the final
ACK, stale passive-handshake reaping, 1024-byte stream segments so UDP
and TURN paths do not rely on IP fragmentation, centralized TURN route
policy, and connection-scoped IPC close reporting for send failures.
The latest Entmoot code consumes those signals through typed stream
classification rather than brittle string matching.

\paragraph{Why the fix is right.} A stream must not disappear without a
reason log, and an accepted stream must already be usable. Once Pilot
made those lifecycle rules explicit, Entmoot's remaining job became
ordinary anti-entropy: drop stale sessions, retry once on a fresh
stream, and fetch the missing bodies. That is why the final deployment
could converge all three stores after the Pilot transport layer went
green.

\section{The big picture: how the layers interacted}
\label{sec:bigpicture}

Many of the bugs documented above \emph{looked like} problems
in Entmoot but were actually problems in Pilot leaking upward.
The four-minute propagation delay we chased on and off for two
days is the best illustration.

Tracing it from the user's perspective back to the wire:

\begin{enumerate}[leftmargin=1.5em]
    \item A user publishes on the VPS. The CLI returns
    instantly (v1.0.3), so the user knows the local store has
    the message, but they can see from the tail on the laptop
    that the message does not arrive.
    \item The Entmoot gossip layer on the VPS runs a fanout. At
    that time the fanout was sequential and had no reachability
    filter, so it tried every peer serially.
    \item Each attempt called Pilot's \texttt{DialConnection}.
    For the laptop peer, the dial landed in a timeout loop
    because (a) the cached UDP address was stale from NAT drift
    on the laptop's carrier, and (b) even when the tunnel was
    \texttt{encrypted=yes auth=yes}, the stream dial timed out
    because Pilot was reply-routing to a stale direct address
    instead of the relay path the last authenticated frame
    arrived on.
    \item Pilot eventually returned a timeout after about 32
    seconds. Entmoot queued a retry. The retry loop redialled
    the same dead peer after a fixed delay, synchronised with
    every other pending message, and paid another 32 seconds.
\end{enumerate}

So the four-minute number was the product of one Pilot bug
(NAT drift refresh, v1.9.0-jf.2; and reply on ingress,
v1.9.0-jf.4), one Entmoot bug (serial fanout with no
reachability filter, v1.0.6), and one Entmoot bug compounded by
a second (deterministic retry schedule without jitter, v1.0.7).
Every patch peeled one layer off the onion. Only after all of
them shipped did the latency number come back to sub-second.

The lesson for the architecture paper is that the split between
layers should be clean enough that operators know where to
look. Most of our time during this shakedown was spent deciding
which layer owned a given symptom. After the dust settled the
rule is simple:

\begin{itemize}[leftmargin=1.2em]
    \item If \texttt{pilotctl peers} shows a peer as
    \texttt{encrypted=yes auth=yes} but \texttt{pilotctl ping}
    still times out, the problem is Pilot.
    \item If \texttt{pilotctl ping} succeeds but Entmoot fanout
    does not deliver, the problem is Entmoot.
\end{itemize}

v1.9.0-jf.4 and v1.0.6 landed in the same day because that is
when we untangled which was which.

The second shakedown refined the rule. If Pilot peers are not
encrypted and authenticated, the failure is still below Entmoot. If
they are authenticated but port~1004 dials time out, first inspect
Pilot stream traces: the bug may be in virtual-stream lifecycle even
though it manifests as an Entmoot reconcile timeout. Only once
\texttt{pilotctl ping} and stream traces are clean should the focus
move to RBSR, store filtering, or gossip retry policy.

\section{Cheat sheet}
\label{sec:cheat}

Two tables, one per layer.

\begin{table}[h]
\centering
\small
\begin{tabular}{@{}llp{7cm}@{}}
\toprule
Version & Date & Summary \\
\midrule
v1.7.2-jf.1 & 18 Apr & Skip NAT punch to private addresses. \\
v1.7.2-jf.2 & 18 Apr & Reset keepalive counter after rekey so tunnels stop flapping every 5 minutes. \\
v1.7.2-jf.3 & 19 Apr & Independent beacon keepalive (25 s); \texttt{set-public} persists to config. \\
v1.8.0-jf.1 & 19 Apr & Pluggable transport layer and TCP fallback. \\
v1.9.0-jf.1 & 19 Apr & Gossip peer-discovery overlay, decoupling reachability from the registry. \\
v1.9.0-jf.2 & 19 Apr & Refresh cached endpoint on every authenticated decrypt; macOS codesign workaround. \\
v1.9.0-jf.3 & 19 Apr & Auto-recovery PILA on crypto-state desync. \\
v1.9.0-jf.4 & 20 Apr & Reply-on-ingress path routing; 3 s SYN retransmit for relay mode. \\
v1.9.0-jf.5 & 20 Apr & Full per-peer state cleanup across \texttt{AddPeer}/\texttt{RemovePeer}. \\
v1.9.0-jf.6 & 21 Apr & \texttt{ErrDialToSelf} guard against self-tunnel storm. \\
v1.9.0-jf.7 & 21 Apr & New \texttt{SetPeerEndpoints} IPC command: Entmoot installs per-peer TCP endpoints into \texttt{peerTCP} without registry mediation. \\
v1.9.0-jf.8--jf.11b & 24--25 Apr & TURN client support, hide-IP posture flags, no-registry endpoint mode, outbound-turn-only semantics, and IPC \texttt{Subscribe}/\texttt{Notify} for TURN rotation. \\
v1.9.0-jf.12--jf.14.2 & 25--26 Apr & Handshake reply routing, keepalives, rendezvous endpoint publish/refresh, and stale endpoint eviction. \\
v1.9.0-jf.15--jf.15.7 & 26 Apr & TURN allocation self-heal, permission refresh diagnostics, bilateral \texttt{CreatePermission}, recovery PILA through TURN, and periodic peer rendezvous re-lookup. \\
v1.9.0-jf.15.8--jf.15.13 & 26--27 Apr & Preserve TURN ingress, close stale IPC stream conns, recover fatal TURN allocation failures, centralize TURN route policy, and prefer own relay in strict TURN-only mode. \\
v1.9.0-jf.15.14--jf.15.21 & 27 Apr & Stream lifecycle tracing, listener \texttt{Unbind}, 1024-byte stream segments, passive-open ACK gating, stale handshake reaping, rendezvous 429 retry, and connection-scoped IPC send failures. \\
\bottomrule
\end{tabular}
\caption{Pilot fork release summary through the \texttt{v1.9.0-jf.15} line.}
\end{table}

\begin{table}[h]
\centering
\small
\begin{tabular}{@{}llp{7cm}@{}}
\toprule
Version & Date & Summary \\
\midrule
v1.0.0 & 17 Apr & SQLite store, control-socket IPC, five-command CLI, invites with TTL. \\
v1.0.2 & 19 Apr & Exponential-backoff retry; anti-entropy reconciliation on reconnect. \\
v1.0.3 & 19 Apr & \texttt{entmootd publish} returns immediately after local store commit. \\
v1.0.4 & 20 Apr & Plumtree dissemination with IHAVE/GRAFT/PRUNE. \\
v1.0.5 & 20 Apr & Re-fanout fires on every acquisition path (push, retry, reconcile). \\
v1.0.6 & 20 Apr & Trust-aware fanout, parallel sends with per-peer timeout, dial-backoff cache. \\
v1.0.7 & 21 Apr & Inline message body on eager push; decorrelated-jitter retry backoff. \\
v1.0.8 & 21 Apr & Exclude issuer from \texttt{bootstrap\_peers} at invite mint. \\
v1.1.0 & 21 Apr & Persistent yamux-multiplexed session per peer; no per-frame Pilot stream dial. \\
v1.2.0 & 21 Apr & Signed \texttt{TransportAd} records gossiped through Entmoot; Pilot \texttt{SetPeerEndpoints} IPC makes TCP fallback actually reachable. \\
v1.2.1 & 22 Apr & Range-Based Set Reconciliation; background AE ticker; \texttt{OnTunnelUp} callback; gap-reconcile bug fixed. \\
v1.3.0 & 25 Apr & Apache-2.0 licensing boundary and independent Pilot IPC client. \\
v1.4.0--v1.4.6 & 25 Apr & Hide-IP transport ads, TURN endpoint publication, cross-layer privacy warnings, TURN polling fallback, and IPC error routing hardening. \\
v1.5.0 & 25 Apr & Pilot \texttt{Subscribe}/\texttt{Notify} support; TURN rotation becomes push-driven with polling only as fallback. \\
v1.5.1--v1.5.4 & 26--27 Apr & Responder-side RBSR body fetch, stale-session recovery, separate dial/I/O deadlines, and post-timeout session eviction. \\
v1.5.5--v1.5.8 & 27 Apr & Preserve early Pilot IPC closes, ignore self transport ads, validate forwarded snapshots, and retry stale Pilot stream IDs without poisoning dial-backoff. \\
v1.5.9--v1.5.10 & 27 Apr & Explicit Pilot/yamux session lifecycle manager, transport trace mode, Pilot readiness wait, reconcile trace mode, and transport-owned stream error classification. \\
\bottomrule
\end{tabular}
\caption{Entmoot release summary through \texttt{v1.5.10}.}
\end{table}

\section*{Appendix A: reading order}

If you only have time for three sections, read
Section~\ref{sec:mental} (the mental model),
Section~\ref{sec:glossary} (the glossary), and
Section~\ref{sec:bigpicture} (how the layers interacted).

If you want to understand what broke on the wire, pick a single
release from Section~\ref{sec:pilot} or~\ref{sec:entmoot}; each
subsection is self-contained and uses the same four-part
template (\emph{what broke}, \emph{what it actually was},
\emph{how we fixed it}, \emph{why the fix is right}).

The CHANGELOGs in \texttt{CHANGELOG.md} and
\texttt{repos/pilotprotocol/CHANGELOG.md} are the source of
truth for the technical claims in this document. The research
paper (\texttt{paper/main.tex}) is the companion to this
companion: the design as intended, with this document as the
field report.

\end{document}
